% Central Paper C content (journal-neutral). Edit here; journal wrappers in ../journals/ only \input these files.
% Every number traces to ../../FINDINGS.md.
\section{Discussion}\label{sec:discussion}

\paragraph{What open data can and cannot show.}
Open citation data reproduce Clarivate's self-citation figures closely when reference lists are present, so open
audits of journal citation behaviour are feasible. But the gaps sit where integrity questions arise: some publishers
deposit no references, and several of the donors behind Clarivate's stacking cases are among them. An open audit can
therefore not conclude that a journal shows no anomaly when its citing partners are absent from the data; coverage
of the relevant journals should be reported alongside any result. Data citations from repositories, such as GBIF
downloads, are a second, newer source of distortion for anyone computing impact-factor-like indicators from open
data.

\paragraph{Why CIDRE misses Clarivate-style stacking.}
CIDRE was designed to find groups whose members exchange excessive citations, and it does find stacking between small
journals. Its symmetric threshold, however, requires a donor to devote a substantial share of its own citations to the
group. In the stacking cases Clarivate reported in 2023, large journals fed small ones, and CIDRE flags the edge but
discards the group. This is a property of the method, not of the data: it appears with well-covered donors, and
lowering the donor threshold merges the pairs into large groups instead of isolating them. A recipient-centric reading
of the same test recovers listed stacking journals far above chance and sometimes before listing, at the cost of
flagging about 2--5\% of journals over nine years, depending on the consensus level. At the same flag rate it
finds more listed journals than a plain rule on the top donor's share (10 against 3 of 28). Part of the raw enrichment
reflects journal size, but five- to seven-fold enrichment remains against size-matched controls. A plain rule with a
volume floor finds 7, not significantly fewer, and flags different journals, so the two are better seen as
complementary screens than as rivals.

\paragraph{What the flags mean.}
The flags mark anomalous citation patterns, not misconduct. Many have benign explanations, such as a society journal
and its companion, editorial preferences, or regional communities. The evidence check shows that flagged pairs carry
independent indicators far more often than matched pairs, and the paper-level anatomy shows that their citations are
spread over many papers that cite the recipient repeatedly and are topically further from it. Flags are candidates for
expert review, not verdicts.

\paragraph{Limitations.}
Clarivate names few stacking pairs, so recall figures rest on a small number of cases, and the enrichment results rely
on 28 to 30 stacking-listed journals. The recipient-centric reading and its consensus rule were chosen after part of
the evaluation. The cases used in the design were excluded and all variants are reported. The prospective check on the 2024
and 2025 lists points the same way but rests on six journals and is not significant; a larger test will become possible
as Clarivate publishes further lists and open data extend beyond 2023. The label table includes nine self-citation rows
from the lists published at the time that Clarivate's current spreadsheet omits; results without them are similar.
The asymmetric-threshold test of CIDRE was evaluated on the 15 named pairs and on journal-level enrichment only.
The evidence check is a lower bound on precision. Clarivate's exact formulas and thresholds are not
public, and the exchange share computed here is an upper bound. Crossref coverage statistics describe the present, not the state at
the time of the snapshot. Title similarity uses TF--IDF on titles only, and the anatomy controls are matched on year and volume, not on
field; the size-matched enrichment uses block membership only as a proxy for field.

